\section{Background}
\label{sec:background}

\textbf{EPSS, KEV, and CVSS.} EPSS, maintained by the Forum of Incident Response and Security
Teams, estimates the probability that a CVE is exploited in the wild~\cite{jacobs2021epss,%
firstepss}. It reframed prioritization away from intrinsic severity and toward predicted
exploitation, and later data-driven work has refined that prediction with richer
features~\cite{jacobs2023enhancing}. The CISA Known Exploited Vulnerabilities
catalog~\cite{cisakev} and Binding Operational Directives 22-01~\cite{cisabod2201} and
23-01~\cite{cisabod2301} mandate timely remediation of confirmed-exploited CVEs and improved
asset visibility, binding agencies to remediation timelines that a fixed workforce must meet
under a finite budget. The Common Vulnerability Scoring System~\cite{cvss31} rates intrinsic
severity independent of any deployment. All three are CVE-level signals: they describe a
vulnerability, not the host it sits on, and so two hosts carrying the same CVE receive the same
exploit-side weight no matter how their mission roles differ. That CVE-level abstraction is the
right one for predicting exploitation, but it is the wrong granularity for deciding which of two
equally exploitable hosts to remediate first when only one can be reached in the window.

\textbf{Hygiene-augmented prioritization.} A hygiene-augmented scorer adds a host-level
hygiene risk score, capturing patch debt, directory exposure, and telemetry freshness, to the
exploit signal~\cite{paper4hygieneprio}. This is the first move from a purely CVE-level view
toward a host-level one, and it measurably raises precision on a hygiene-defined target by
surfacing the poorly-maintained hosts that an exploit-only ranking would miss. But the host
weight it introduces encodes the probability that a host is a weak link, not the consequence of
its compromise: a host scores high because it is badly patched and exposed, not because it
matters. A poorly-patched training kiosk can therefore outrank a well-patched CJIS controller,
which inverts the order an authorizing official would choose. CAP-G keeps this hygiene-augmented
scorer as its base and is measured strictly as the increment that asset context adds on top of
it; the base scorer is defined self-containedly in Section~\ref{sec:method} so that the paper is
readable without its predecessor~\cite{paper4hygieneprio,paper1vulnprio}.

\textbf{Asset criticality in practice.} FIPS~199~\cite{fips199} and SP~800-60~\cite{nist80060}
categorize systems by impact level across confidentiality, integrity, and availability, and the
federal Risk Management Framework built on SP~800-53~\cite{nist80053} requires that remediation
be driven by that impact categorization rather than by raw severity. The CJIS Security
Policy~\cite{cjis2023} imposes elevated controls on criminal-justice systems, making the
data-sensitivity dimension of context a compliance fact rather than a modelling convenience. The
attributes this paper uses, asset criticality tier, network exposure zone, and data sensitivity,
are precisely the ones public-sector configuration management databases already record for inventory
and authorization purposes. CAP-G therefore operationalizes already-collected attributes as a
ranking signal rather than introducing new instrumentation: the cost of adopting it is a scoring
change, not a data-collection program. Reusing an existing signal in this way also keeps the
intervention auditable, since every input to the ranking is an attribute the organization already
maintains and can explain.
